% ---------------------------------------------------------------
% preambulo.tex — pacotes, tipos de letra, cores e estilos
% Tipos de letra iguais aos do Word: Book Antiqua (texto),
% Georgia (títulos) e Aptos (legendas). Requer XeLaTeX/Tectonic.
% ---------------------------------------------------------------

\usepackage{fontspec}
\usepackage[portuguese]{babel}
\usepackage{microtype}

\setmainfont{Book Antiqua}[AutoFakeBold=1.5,AutoFakeSlant=0.2]
\setmonofont{Consolas}[Scale=0.9]
\newfontfamily\titulos{Georgia}
\newfontfamily\legendas{Aptos}[Path=fontes/,Extension=.ttf,
  UprightFont=*-Regular,BoldFont=*-Bold,ItalicFont=*-Italic,BoldItalicFont=*-BoldItalic]
% Símbolos que a Book Antiqua não tem (setas, etc.)
\newfontfamily\simbolosB{Cambria}
\newfontfamily\simbolos{Segoe UI Symbol}

\usepackage[a4paper,margin=2.5cm]{geometry}
\usepackage{graphicx}
\usepackage{array}
\usepackage{longtable}
\usepackage[table]{xcolor}
\usepackage{float}
\usepackage{pdflscape}
\usepackage{enumitem}
\usepackage{caption}
\usepackage{titlesec}
\usepackage[titles]{tocloft}
\usepackage{fancyhdr}
\usepackage{csquotes}
\usepackage{url}
\usepackage[most]{tcolorbox}
\usetikzlibrary{positioning,arrows.meta,calc,fit}

% --- Cores -----------------------------------------------------
\definecolor{corEscura}{HTML}{1F3A5F}
\definecolor{corDestaque}{HTML}{AD451D}
\definecolor{corCinza}{HTML}{5A6672}
\definecolor{corFundo}{HTML}{F4F6F8}
\definecolor{corLegenda}{HTML}{404040}

% --- Parágrafos: texto justificado, sem avanço, espaço entre parágrafos
\setlength{\parindent}{0pt}
\setlength{\parskip}{6pt plus 2pt}
\linespread{1.08}
\setlength{\emergencystretch}{2em}
\setlist{itemsep=2pt,parsep=0pt,topsep=4pt}

% --- Títulos (Georgia) -----------------------------------------
\titleformat{\chapter}[hang]
  {\titulos\fontsize{20}{24}\selectfont\bfseries\color{corEscura}}{\thechapter.}{0.5em}{}
\titleformat{name=\chapter,numberless}[hang]
  {\titulos\fontsize{20}{24}\selectfont\bfseries\color{corEscura}}{}{0pt}{}
\titlespacing*{\chapter}{0pt}{0pt}{18pt}
\titleformat{\section}
  {\titulos\fontsize{12.5}{15}\selectfont\bfseries\color{corEscura}}{\thesection.}{0.5em}{}
\titlespacing*{\section}{0pt}{16pt}{6pt}
\titleformat{\subsection}
  {\titulos\fontsize{11.5}{14}\selectfont\bfseries\itshape\color{corEscura}}{\thesubsection.}{0.5em}{}

% Anexo: \anexo{A}{Título} — título sem número de capítulo, entrada no índice.
\newcommand{\anexo}[2]{\chapter*{Anexo #1 – #2}\addcontentsline{toc}{chapter}{Anexo #1 – #2}}

% --- Legendas (Aptos, centradas) -------------------------------
\DeclareCaptionFont{legenda}{\legendas\fontsize{9}{11}\selectfont\color{corLegenda}}
\captionsetup{font=legenda,labelsep=endash,justification=centering,singlelinecheck=false,skip=6pt}
\captionsetup[table]{position=top}

% Figuras e tabelas numeradas de forma seguida (como no Word).
\counterwithout{figure}{chapter}
\counterwithout{table}{chapter}

% --- Índice e listas (como no Word) ----------------------------
\addto\captionsportuguese{%
  \renewcommand{\contentsname}{Índice}%
  \renewcommand{\listfigurename}{Lista de figuras}%
  \renewcommand{\listtablename}{Lista de tabelas}%
  \renewcommand{\bibname}{Bibliografia}%
}
\renewcommand{\cfttoctitlefont}{\titulos\fontsize{20}{24}\selectfont\bfseries\color{corEscura}}
\renewcommand{\cftloftitlefont}{\cfttoctitlefont}
\renewcommand{\cftlottitlefont}{\cfttoctitlefont}
\renewcommand{\cftchapdotsep}{\cftdotsep}
\renewcommand{\cftfigpresnum}{Figura }
\renewcommand{\cftfigaftersnum}{ –}
\setlength{\cftfignumwidth}{5.6em}
\renewcommand{\cfttabpresnum}{Tabela }
\renewcommand{\cfttabaftersnum}{ –}
\setlength{\cfttabnumwidth}{5.6em}
\setlength{\cftfigindent}{0pt}
\setlength{\cfttabindent}{0pt}
% Sem espaço entre capítulos nas listas de figuras e de tabelas (como no Word).
\usepackage{etoolbox}
\makeatletter
\patchcmd{\@chapter}{\addtocontents{lof}{\protect\addvspace{10\p@}}}{}{}{}
\patchcmd{\@chapter}{\addtocontents{lot}{\protect\addvspace{10\p@}}}{}{}{}
\makeatother
\urlstyle{same}

% --- Número de página centrado no rodapé -----------------------
\pagestyle{plain}

% --- Tabelas (estilo do Word: cabeçalho azul, linhas alternadas)
\newcolumntype{L}[1]{>{\raggedright\arraybackslash}p{#1}}
\newcommand{\cab}[1]{{\color{white}\bfseries #1}}
\newcommand{\linhacab}{\rowcolor{corEscura}}
% \begin{tabela}{colunas} ... \end{tabela}: longtable em letra pequena, com quebra de página.
\newenvironment{tabela}[1]
  {\par\begingroup\footnotesize\renewcommand{\arraystretch}{1.35}%
   \setlength{\LTpre}{6pt}\setlength{\LTpost}{6pt}%
   \arrayrulecolor{corEscura}\begin{longtable}{#1}}
  {\end{longtable}\endgroup}

% --- Caixas e macros -------------------------------------------
% Estrutura das regras (SE ... ENTÃO ...).
\newtcolorbox{caixaregra}{enhanced, colback=corFundo, colframe=corEscura,
  boxrule=0pt, leftrule=1.2pt, arc=0pt, outer arc=0pt, left=8pt, right=8pt,
  top=4pt, bottom=4pt, fontupper=\ttfamily\small, width=0.75\linewidth, center}
% Versaletes simuladas (a Book Antiqua não as tem), como faz o Word: \pv{texto} escreve
% minúsculas como maiúsculas pequenas (80%). Usado na capa.
\newcommand{\pv}[1]{{\fontsize{0.8\dimexpr\csname f@size\endcsname pt\relax}{0}\selectfont #1}}
% Nomes de componentes e identificadores do código.
\newcommand{\cod}[1]{\texttt{#1}}

\usepackage[hidelinks,bookmarks=true]{hyperref}
\hypersetup{pdftitle={\tituloprincipal},pdfauthor={\autoresmeta},pdflang={pt-PT}}
